\clearpage
\section*{ВИСНОВКИ}
\phantomsection
\addcontentsline{toc}{section}{ВИСНОВКИ}

У ході лабораторної роботи досліджено особливості розподілених транзакцій і властивості ACID у середовищі з кількома вузлами. Порівняльний аналіз показав, що двофазний протокол фіксації є придатним для моделювання банківського переказу, оскільки забезпечує єдине рішення для всіх учасників операції. Водночас його використання потребує контролю часу перебування транзакцій у стані \texttt{PREPARED}, оскільки підготовлені операції можуть утримувати блокування.

Для PostgreSQL сформовано модель із двох баз даних і координатора. Успішний сценарій завершує зміну балансів через \texttt{COMMIT PREPARED}: сума \(250.00~\mathrm{UAH}\) списується з рахунку відправника та зараховується на рахунок отримувача. У сценарії відмови команда \texttt{ROLLBACK PREPARED} скасовує підготовлені зміни, унаслідок чого часткове виконання переказу не виникає.

Застосування журналу рішень координатора дає змогу відновити завершення транзакції після збою зв'язку або перезапуску програми. Отже, поставлену мету досягнуто: принципи 2PC проаналізовано, модель розподіленого переказу побудовано, а очікувані результати успішної фіксації та відкату визначено.